% SPDX-License-Identifier: CC-BY-4.0
\section{Theory: Fisher information, the datum defect, and rank restoration}
\label{sec:theory}

The empirical contrast in the results section---absolute station error collapsing
by orders of magnitude once Earth-baseline VLBI is switched on for a sparse
constellation---asks for a reason. It is not that VLBI brings more numbers to the
fit. It is that without it the fit is missing a particular piece of information,
and the information geometry says exactly which piece, exactly when it is missing,
and exactly how much is needed to supply it. This section makes that precise. The
statements here are exact algebra of the estimation problem, and each is checked
against the engine that produces the figures: where we claim a subspace lies in
the null space of the information matrix, the engine's computed null space confirms
it to numerical precision.

\subsection{The Fisher information of the joint solve}

Write the estimated state as the stacked correction vector
$x = [\,\delta r_S,\ \{\delta r_k\},\ c\tau_S,\ \{c\tau_k\}\,]$ (station position,
satellite positions, station and satellite clocks carried as range-equivalent
metres), and the observations as $y = h(x) + n$ with $n \sim \mathcal{N}(0,R)$ and
$R = \operatorname{diag}(\sigma_i^2)$. For this Gaussian model the Fisher
information matrix is
\begin{equation}
  M = H^{\top} W H, \qquad W = R^{-1}, \qquad H = \frac{\partial h}{\partial x},
  \label{eq:fim}
\end{equation}
and the Cram\'er--Rao lower bound states that any unbiased estimator has
covariance $\operatorname{Cov}(\hat x) \succeq M^{-1}$. When $M$ is rank-deficient
the bound on the observable subspace is the Moore--Penrose pseudo-inverse $M^{+}$,
and the unobservable directions---the null space of $M$---have no finite bound at
all. The question ``is the station's absolute position observable?'' is therefore
literally ``does the station-position block of $x$ lie outside $\ker M$?'' Kshana
computes $M$, its spectrum, $M^{+}$ and the null-space basis directly
(\code{src/fim.rs}); that engine is Validated against NumPy's \code{eigh}/\code{inv}
and the published closed-form bounds of \citet{kay1993} (\S\ref{sec:toolkit}).

The information is additive across techniques, $M = \sum_i w_i\, h_i' h_i'^{\top}$,
so each observable contributes a rank-one outer product of its state gradient. Five
kinds of observable feed it: a station clock-sync pseudo-observable; lunar-local
ranges (station$\leftrightarrow$satellite and inter-satellite); Earth$\to$satellite
radiometric ranges; and the Earth-baseline VLBI delays to the station beacon. The
argument below tracks which datum directions each kind can and cannot constrain.

\subsection{The datum defect is the rigid-body group}

Start with the lunar-local ranges alone. A range $\lvert r_a - r_b\rvert$ between
two members of the lunar cluster (station or satellites) depends only on their
relative geometry. Translate the whole cluster rigidly, or rotate it rigidly about
the Moon's centre, and not one station$\leftrightarrow$satellite or inter-satellite
range changes. The rigid motions of a three-dimensional point cluster form a
six-parameter group---three translations and three rotations---so the lunar-local
information has a null space of dimension (at least) six. This is the lunar
instance of the free-network, or rank-deficiency, problem of geodetic adjustment:
internal observations fix internal geometry and leave the absolute datum free.

We state this as the load-bearing fact of the paper and check it directly. Let
$H_{\mathrm{int}}$ be the Jacobian of the lunar-local observables (station$\to$%
satellite plus inter-satellite ranges) only, and let
$g_1,\dots,g_6$ be the six rigid-body generators: the three unit translations of
every cluster point, and the three infinitesimal rotations $g = e_a \times (r_p -
\bar r)$ about the cluster centroid $\bar r$.

\begin{theorem}[Internal datum defect]
\label{thm:rigid}
For a non-degenerate lunar cluster the six rigid-body generators are linearly
independent and satisfy $H_{\mathrm{int}}\, g_m = 0$ for $m = 1,\dots,6$. They
therefore span a six-dimensional subspace of $\ker M$, and because the rigid-body
group is the only continuous family of maps that preserves every pairwise range,
it is the complete positional datum defect of the lunar-local geometry.
\end{theorem}

The translations are exactly invariant by construction; the rotations are invariant
to first order, which is exact for the information matrix. Computing
$H_{\mathrm{int}}$ for the six-satellite network and applying it to the six
generators, the engine returns a relative residual of $0$ for each translation and
$\sim\!10^{-9}$ for each rotation (the finite-difference floor of the Jacobian),
while a random non-rigid direction returns a residual of order unity. The datum
defect is the rigid-body group, not an artefact and not a coincidence of
dimension.

A second observation removes the clocks from the discussion. Ranges measure only
clock \emph{differences}, so a common offset added to every clock is unobservable;
the single station clock-sync pseudo-observable pins that one direction. With the
sync in place every clock is observable, and the residual datum defect lives
\emph{entirely in position}. The engine confirms this: the null-space basis carries
zero energy in the clock coordinates. From here ``datum defect'' means the
unobservable \emph{positional} subspace.

\subsection{Closing the datum: two Earth-frame ties}

Nothing internal to the Moon can fix the absolute datum; it has to be tied to the
Earth/inertial frame. Two observables do that, and they are not equivalent.

The \emph{indirect} tie is Earth$\to$satellite ranging. It anchors the
\emph{satellites} in the geocentric frame, and lunar-local ranging then propagates
that anchor to the station. Whether it reaches the station's absolute position
depends on the constellation geometry: the propagation is only as strong as the
station$\to$satellite look-direction spread. For a sparse constellation confined to
one part of the station's sky the propagation is rank-deficient---one station
direction never gets tied---and the station's absolute position stays in the null
space no matter how precise the ranges are.

The \emph{direct} tie is the Earth-baseline VLBI delay to the station beacon. Each
delay is the difference of two near-parallel look directions from the lunar beacon
to two Earth stations, so it constrains the station's absolute position itself,
independent of the constellation. Each snapshot baseline contributes a single
rank-one term; a delay \emph{rate} (an Earth-rotation arc, not the single epoch
modelled here) would add a second. Because the Earth stations subtend a tiny angle
from the Moon, each baseline gradient lies in the plane of the sky, so one epoch of
VLBI informs at most two of the station's three absolute axes, and at least two
non-collinear baselines---three non-collinear Earth stations---are needed to span
the residual datum.

The consequence is a clean, testable claim: \emph{VLBI restores observability when
the indirect tie cannot, and merely sharpens it when the indirect tie already can.}
Table~\ref{tab:restore} is the engine's verdict, and it is exactly this.

\begin{table}[t]
\centering
\small
\begin{tabular}{llcc}
\toprule
Constellation & Earth stations & Without VLBI & With VLBI \\
\midrule
Sparse (3 sat) & 6 & \textbf{unobservable} (defect 1) & observable, \SI{20.1}{\meter} \\
Rich (6 sat)   & 6 & observable, \SI{23.2}{\meter}   & observable, \SI{9.7}{\meter} \\
\bottomrule
\end{tabular}
\caption{Whether the station's absolute 3-D position is observable, from the
Fisher information (Cram\'er--Rao bound on the absolute position, $1\sigma$ RSS).
For a sparse constellation VLBI \emph{restores} observability---without it the
position lies in $\ker M$ and has no finite bound. For a rich constellation the
indirect tie already makes it observable and VLBI \emph{sharpens} it
($\sim\!2.4\times$ in the bound). The design question is therefore not ``is VLBI
nice to have'' but ``is the constellation rich enough to do without it.''}
\label{tab:restore}
\end{table}

\subsection{The three-station threshold and the design curve}

Take the regime where VLBI carries the absolute datum---a sparse or early
constellation, or a conservative design that does not lean on rich satellite
geometry. Sweeping the number of Earth stations $n_E$ with the direct tie active,
and using a network with enough satellites that observation count is not the
limiting factor, the positional datum defect falls
\begin{equation}
  n_E = 1:\ \text{defect } 3 \;\longrightarrow\;
  n_E = 2:\ \text{defect } 1 \;\longrightarrow\;
  n_E = 3:\ \text{defect } 0,
  \label{eq:defect}
\end{equation}
so the absolute station position becomes observable at three non-collinear Earth
stations and not before. Each new non-collinear station contributes plane-of-sky
baseline directions that close part of the residual datum, with the last direction
closing at the third station---the parallax content of the geometry, read straight
off $\dim\ker M$.

One caveat keeps the ladder honest. A network with too few satellites is also
\emph{observation-starved}: it has fewer measurements than parameters, which
inflates the apparent defect. For the three-satellite default the \emph{observed}
defect runs $6 \to 2 \to 0$, but that is the true datum defect of \eqref{eq:defect}
plus an equal observation-count deficit ($3 \to 1 \to 0$) that has nothing to do
with the datum. We report \eqref{eq:defect}, the true datum defect, measured on a
network rich enough that the two effects separate; the engine reproduces the
decomposition exactly. The three-station threshold itself is robust---it holds at
three, six and eight satellites alike.

Once the position is observable the bound sets the achievable accuracy. For the
six-satellite network the Cram\'er--Rao bound on the station's absolute 3-D position
falls from \SI{12.1}{\meter} at three Earth stations to \SI{9.7}{\meter} at six and
\SI{8.5}{\meter} at ten---the diminishing-returns design curve of
Figure~\ref{fig:n_earth}, and the basis of the sizing rule in \S\ref{sec:discussion}.

\subsection{The bound is attained}

The bound is not a floor the estimator sits comfortably above. For the
six-satellite, six-station network the Cram\'er--Rao bound on the absolute station
position is \SI{9.7}{\meter} ($1\sigma$ RSS), the Gauss--Newton estimator's
Monte-Carlo RMS over 200 seeds is \SI{9.9}{\meter}, an efficiency
$\text{RMS}/\text{CRLB} = 1.02$, and the empirical bias is \SI{0.36}{\meter}---%
under four percent of the RMS, so the estimator is effectively unbiased and the
RMS-to-CRLB comparison is the right one. The estimator is statistically efficient:
it extracts essentially all the information the geometry holds, so the gains
reported next are the information-theoretic optimum for each geometry, not an
artefact of a particular solver.

One honesty note carries through the rest of the paper. Everything above is exact
algebra of the estimation model, checked against the engine that computes it; it
explains the simulated result, it does not validate it against real lunar tracking.
The FIM/CRLB engine is Validated against independent libraries, but the lunar
network it is applied to is a representative simulation, so the observability result
itself stays Modelled. We keep that line bright in \S\ref{sec:toolkit}.
